% 子文件：附录.tex
% 使用 \input{附录.tex} 引入，无 document 环境

\section{Stolz定理}

\subsection{定理陈述}
设$\{a_n\}$和$\{b_n\}$是两个数列，且满足以下条件：
\begin{enumerate}
    \item 数列$\{b_n\}$单调增加并趋于无穷大；
    \item 存在极限 $\displaystyle\lim_{n\rightarrow\infty} \frac{a_{n+1}-a_n}{b_{n+1}-b_n} = L$。
\end{enumerate}
那么，若$\{b_n\}$的增量$b_{n+1}-b_n$满足$\displaystyle\sum_{n=1}^{\infty}(b_{n+1}-b_n)=\infty$，则有
\[
\lim_{n\rightarrow\infty}\frac{a_n}{b_n}=L.
\]

\subsection{证明思路}
Stolz定理的证明通常基于柯西收敛原理和单调有界原理。
\begin{enumerate}
    \item 假设前提成立：设有数列$\{a_n\}$和$\{b_n\}$，$\{b_n\}$单调递增且趋于无穷大，且
    \[
    \lim_{n\rightarrow\infty}\frac{a_{n+1}-a_n}{b_{n+1}-b_n}=L.
    \]
    \item 构造辅助数列：定义 $\displaystyle d_n=\frac{a_n}{b_n}$。
    \item 应用柯西收敛原理：$\forall \varepsilon>0$，$\exists N\in\mathbb{N}^*$，当$n>N$时，
    \[
    \left|\frac{a_{n+1}-a_n}{b_{n+1}-b_n}-L\right|<\varepsilon.
    \]
    \item 估计差异：对任意$m>n>N$，$b_m-b_n>0$，于是
    \[
    \left|\frac{a_m}{b_m}-\frac{a_n}{b_n}\right|
    =\left|\frac{a_mb_n -a_nb_m}{b_mb_n}\right|
    \le\left|\frac{(a_{n+1}-a_n)+(a_{n+2}-a_{n+1})+\cdots+(a_m-a_{m-1})}{b_m-b_n}\right|.
    \]
    \item 应用极限定义：
    \[
    \left|\frac{(a_{n+1}-a_n)+(a_{n+2}-a_{n+1})+\cdots+(a_m-a_{m-1})}{b_m-b_n}-L\right|<\varepsilon.
    \]
    \item 结论：$\displaystyle\lim_{n\rightarrow\infty}\frac{a_n}{b_n}=L$。
\end{enumerate}

\section{柯西收敛原理}
柯西收敛原理是实数系中判断数列收敛的充要条件。

设$\{a_n\}_{n=1}^{\infty}$是一个数列，如果$\forall \varepsilon>0$，存在正整数$N$，使得对一切$m,n>N$，都有
\[
|a_m - a_n|<\varepsilon,
\]
则称$\{a_n\}$为柯西数列。

\textbf{柯西收敛原理：} 在实数域上，数列$\{a_n\}$收敛的充分必要条件是$\{a_n\}$为柯西数列。

\section{达布定理(Darboux定理)}
若函数$f$在闭区间$[a,b]$上连续，在开区间$(a,b)$内可导，则对介于$f'(a),f'(b)$之间的任意实数$m$，至少存在一点$c\in(a,b)$，使得
\[
f'(c)=m.
\]
>说明：导函数具有介值性，导函数不会发生第一类跳跃间断。

\section{常用积分公式}
\subsection{区间再现公式}
\[
\int_a^b f(x)\,dx=\int_a^b f(a+b-x)\,dx
\]

\subsection{华里士公式（点火公式）}
\[
\int_{0}^{\frac{\pi}{2}} \sin^n x \, dx = \int_{0}^{\frac{\pi}{2}} \cos^n x \, dx
=
\begin{cases}
\displaystyle \frac{(n-1)!!}{n!!} \cdot \frac{\pi}{2}, & n\text{ 为偶数} \\[6pt]
\displaystyle \frac{(n-1)!!}{n!!}, & n\text{ 为奇数}
\end{cases}
\]

推广：
\[
\int_{0}^{\frac{\pi}{2}} \sin^n x \cos^m x \, dx
=
\begin{cases}
\displaystyle\frac{(m-1)!!\,(n-1)!!}{(m+n)!!}\cdot \frac{\pi}{2}, & m,n\text{ 全为偶数}\\[6pt]
\displaystyle\frac{(m-1)!!\,(n-1)!!}{(m+n)!!}, & m,n\text{ 不全为偶数}
\end{cases}
\]

\subsection{表格积分法（分部积分推广）}
设$u(x),v(x)$均有直到$n+1$阶导数，则
\[
\int u\, v^{(n+1)} dx
\]
\[
\begin{array}{c|c}
\text{求导}u & \text{积分}v^{(n+1)} \\
\hline
u & v^{(n+1)} \\
u' & v^{(n)} \\
u'' & v^{(n-1)} \\
\vdots & \vdots
\end{array}
\]

\section{欧拉公式}
\[
e^{i\theta}=\cos\theta+i\sin\theta,\qquad
e^{-i\theta}=\cos\theta-i\sin\theta
\]

\section{重积分的变量代换}
% p87 笔记占位
设变换 $T:\begin{cases}x=x(u,v)\\y=y(u,v)\end{cases}$，雅可比行列式
\[
J=\frac{\partial(x,y)}{\partial(u,v)}
=\begin{vmatrix}
\dfrac{\partial x}{\partial u} & \dfrac{\partial x}{\partial v}\\[4pt]
\dfrac{\partial y}{\partial u} & \dfrac{\partial y}{\partial v}
\end{vmatrix}
\]
\[
\iint_D f(x,y)\,dxdy
=\iint_{D'} f\big(x(u,v),y(u,v)\big)\,|J|\,dudv
\]

\section{含参变量的积分}
设函数
\[
I(y)=\int_a^b f(x,y)\,dx
\]

\section{常见幂级数展开式}
\[
\ln(1+x)=\sum_{n=1}^{\infty}\frac{(-1)^{n-1}}{n}x^n,\quad |x|<1
\]
\[
e^x=\sum_{n=0}^{\infty}\frac{x^n}{n!},\quad x\in\mathbb R
\]
\[
\cos x=\sum_{n=0}^{\infty}\frac{(-1)^n}{(2n)!}x^{2n},\quad x\in\mathbb R
\]
\[
\frac{1}{1+x}=\sum_{n=0}^{\infty} (-1)^n x^n,\quad |x|<1
\]